% AUTORIDAD DE FUENTE: Exposé IX, §7.3, líneas 1969--2015;
% índices de escaneo 418--420, folios 407--409. Las capas francesas D/C
% no presentan diferencias materiales.

\medskip
\noindent
7.3. Consideremos ahora, para un número primo fijo $\ell$, el grupo
pro-$\ell$ $T_{\ell}(A^{o})$. Identificando un esquema finito sobre $S$
con el esquema formal correspondiente, y definiendo como de costumbre
la ``\underline{parte fija}'' $T_{\ell}(A^{o})^{f}$, obtenemos un
isomorfismo como en (5.1.5):

%%%%%%%%%%  scan idx 419  |  folio 408  %%%%%%%%%%

\begin{equation}\tag{7.3.1}
T_{\ell}(A^{o})^{f} \simeq T_{\ell}(\widehat{A}^{o}) \quad ,
\end{equation}
y, definiendo la ``\underline{parte tórica}'' como en (5.1.6),

\begin{equation}\tag{7.3.2}
T_{\ell}(A^{o})^{t} = T_{\ell}(\widehat{T}) \simeq T_{\ell}(T)
 \subset T_{\ell}(A^{o})^{f} \simeq T_{\ell}(\widehat{A}^{o}) \quad ,
\end{equation}
obtenemos inmediatamente, para la ``\underline{parte abeliana}'', un
isomorfismo

\begin{equation}\tag{7.3.3}
T_{\ell}(A^{o})^{ab} \stackrel{\mathrm{def}}{=}
T_{\ell}(A^{o})^{f}/T_{\ell}(A^{o})^{t}
 \simeq T_{\ell}(\widehat{B}) \quad .
\end{equation}
Hasta aquí aún no hemos utilizado la hipótesis de que $\widehat B$ sea
algebraizable, hipótesis que nos permitió construir el esquema abeliano
$B$ y la extensión de Raynaud de $B$ por $T$. Puesto que existe,
naturalmente, un isomorfismo canónico (el caso particular de (7.3.1)
aplicado al esquema abeliano $B$)
$T_{\ell}(\widehat B)\simeq T_{\ell}(B)$, el isomorfismo (7.3.3) también
nos da

\begin{equation}\tag{7.3.4}
T_{\ell}(A^{o})^{ab} \simeq T_{\ell}(B) \quad .
\end{equation}
Señalemos asimismo, de paso, que (7.2.3) y (7.3.1) implican un
isomorfismo canónico

\begin{equation}\tag{7.3.5}
T_{\ell}(A^{o})^{f} \simeq T_{\ell}(A^{\natural o}) \quad ,
\end{equation}
pues el miembro de la derecha ya es un grupo profinito, es decir,
idéntico a su parte fija. La sucesión exacta canónica

\begin{equation}\tag{7.3.6}
0 \longrightarrow T_{\ell}(A^{o})^{t}
 \longrightarrow T_{\ell}(A^{o})^{f}
 \longrightarrow T_{\ell}(B) \longrightarrow 0 \quad ,
\end{equation}
que, en cualquier caso (sin la hipótesis de algebraizabilidad de
$\widehat B$), puede interpretarse como la sucesión exacta obtenida al
aplicar el funtor $T_{\ell}$ a la sucesión exacta

\begin{equation*}
0 \longrightarrow \widehat{T} \longrightarrow \widehat{A}^{o}
 \longrightarrow \widehat{B} \longrightarrow 0 \quad ,
\end{equation*}

%%%%%%%%%%  scan idx 420  |  folio 409  %%%%%%%%%%

también puede considerarse obtenida al aplicar el funtor $T_{\ell}$ a
la sucesión exacta canónica de esquemas en grupos ordinarios

\begin{equation*}
0 \longrightarrow T \longrightarrow A^{\natural o}
 \longrightarrow B \longrightarrow 0 \qquad .
\end{equation*}
